\documentclass[../../main.tex]{subfiles}
\graphicspath{{../../image/}{../../figure/}}
\begin{document}

\section{等距浸入}\label{section:ch18-等距浸入}

微分几何是混合型椭圆--双曲平衡律方程组的另一个丰富源泉。此处我们考虑经典问题：将二维 Riemann 流形等距嵌入到 $\mathbb{R}^3$ 中。

考虑二维流形 $\Omega\subset\mathbb{R}^2$，配以光滑 Riemann 度量 $[g_{ij}]$。$[g_{ij}]$ 的逆矩阵记为 $[g^{kl}]$，$|g|$ 表示 $\det[g_{ij}]$。我们采用微分几何的标准记号，包括求和约定。

度量诱导出 Christoffel 符号
\begin{align}
\Gamma_{ij}^k=\frac12 g^{kl}\bigl(\partial_j g_{il}+\partial_i g_{jl}-\partial_l g_{ij}\bigr),
\label{align:ch18-1871}
\end{align}
进而得曲率张量
\begin{align}
R_{ijkl}=g_{lm}\bigl(\partial_k\Gamma_{ij}^m-\partial_j\Gamma_{ik}^m+\Gamma_{ij}^n\Gamma_{nk}^m-\Gamma_{ik}^n\Gamma_{nj}^m\bigr)
\label{align:ch18-1872}
\end{align}
以及 Gauss 曲率
\begin{align}
\kappa=\frac{R_{1212}}{|g|}.
\label{align:ch18-1873}
\end{align}

流形到 $\mathbb{R}^3$ 的浸入 $r$ 保持第一基本形式 $I=g_{ij}\,\mathrm{d}x^i\mathrm{d}x^j$，即
\begin{align}
\partial_i r\cdot\partial_j r=g_{ij},\qquad i,j=1,2,
\label{align:ch18-1874}
\end{align}
当且仅当第二基本形式 $II=h_{ij}\,\mathrm{d}x^i\mathrm{d}x^j$ 满足经典 Gauss--Codazzi 方程。为方便计，我们把 $(x^1,x^2)$ 重新记为 $(x,y)$，并将 $h_{ij}$ 重新标度为
\begin{align}
L=\frac{h_{11}}{\sqrt{|g|}},\qquad M=\frac{h_{12}}{\sqrt{|g|}},\qquad N=\frac{h_{22}}{\sqrt{|g|}},
\label{align:ch18-1875}
\end{align}
此时 Gauss--Codazzi 方程取形式
\begin{align}
\begin{cases}
\partial_x N-\partial_y M=-\Gamma_{22}^1 L+2\Gamma_{12}^1 M-\Gamma_{11}^1 N\\
\partial_x M-\partial_y L=\Gamma_{22}^2 L-2\Gamma_{12}^2 M+\Gamma_{11}^2 N,
\end{cases}
\label{align:ch18-1876}
\end{align}
\begin{align}
LN-M^2=\kappa.
\label{align:ch18-1877}
\end{align}

视相伴的 $(L,M,N)$ 是方程组\eqref{align:ch18-1876}、\eqref{align:ch18-1877} 的 Lipschitz 连续经典解还是仅本性有界可测弱解，浸入称为正则的或奇异的。此处我们在奇异浸入的范围内操作。

一个关键观察是：弱解序列的 $L^\infty$ 弱$^*$ 极限本身仍是弱解。这由 div--curl 引理（\cref{theorem:ch17-divcurl引理}）得到：只需把\eqref{align:ch18-1876}$_1$、\eqref{align:ch18-1876}$_2$ 与\eqref{align:ch18-1877} 的左端分别实现为 $\operatorname{div}(N,-M)$、$\operatorname{curl}(L,M)$ 与 $(N,-M)\cdot(L,M)$。这就为用补偿紧致方法处理该方程组铺平了道路。

从守恒律理论的立场看，可把\eqref{align:ch18-1876}实现为一个平衡律方程组，它可通过为 $L,M,N$ 指定“本构方程”并辅以相容约束\eqref{align:ch18-1877} 而封闭。为此，我们选取状态向量 $(\rho,u,v)$ 以及一个单调函数 $p(\rho)$，令
\begin{align}
L=\rho v^2+p(\rho),\qquad M=-\rho uv,\qquad N=\rho u^2+p(\rho),
\label{align:ch18-1878}
\end{align}
此时\eqref{align:ch18-1876}取形式
\begin{align}
\begin{cases}
\partial_x(\rho u^2+p(\rho))+\partial_y(\rho uv)=X\\
\partial_x(\rho uv)+\partial_y(\rho v^2+p(\rho))=Y,
\end{cases}
\label{align:ch18-1879}
\end{align}
其中
\begin{align}
\begin{cases}
X=-(\rho v^2+p)\Gamma_{22}^1-2\rho uv\Gamma_{12}^1-(\rho u^2+p)\Gamma_{11}^1\\
Y=-(\rho v^2+p)\Gamma_{22}^2-2\rho uv\Gamma_{12}^2-(\rho u^2+p)\Gamma_{11}^2.
\end{cases}
\label{align:ch18-18710}
\end{align}

于是\eqref{align:ch18-1876}退化为状态方程为 $p(\rho)$ 的等熵定常无旋平面气流的动量平衡方程。事实上，\eqref{align:ch18-1879} 与\eqref{align:ch18-1821}$_{2,3}$ 类似，$(X,Y)$ 扮演体积力的角色。

状态方程取为
\begin{align}
p(\rho)=-\frac1\rho,
\label{align:ch18-18711}
\end{align}
它表示一种 Chaplygin 型气体\eqref{align:ch02-2-5-23}。为满足\eqref{align:ch18-1877}，由\eqref{align:ch18-1878}与\eqref{align:ch18-18711}得
\begin{align}
-q^2+\frac{1}{\rho^2}=\kappa,\qquad q^2=u^2+v^2.
\label{align:ch18-18712}
\end{align}
状态方程为\eqref{align:ch18-18711}的气体焓为 $-\frac{1}{2\rho^2}$，故\eqref{align:ch18-18712}可视为 Bernoulli 方程，与\eqref{align:ch18-18214}类似，只是多了一个“体积力” $\frac12\operatorname{grad}\kappa$。

为完成等距浸入与气流之间的类比，我们导出由\eqref{align:ch18-1879}与\eqref{align:ch18-18712}诱导的质量与涡量平衡律。联立这些方程得
\begin{align}
\begin{cases}
u[\partial_x(\rho u)+\partial_y(\rho v)]-\rho v[\partial_x v-\partial_y u]=\frac12\rho\partial_x\kappa+X\\
v[\partial_x(\rho u)+\partial_y(\rho v)]+\rho u[\partial_x v-\partial_y u]=\frac12\rho\partial_y\kappa+Y,
\end{cases}
\label{align:ch18-18713}
\end{align}
从而
\begin{align}
\partial_x(\rho u)+\partial_y(\rho v)=\frac{\rho u}{2q^2}\partial_x\kappa+\frac{\rho v}{2q^2}\partial_y\kappa+\frac{u}{q^2}X+\frac{v}{q^2}Y,
\label{align:ch18-18714}
\end{align}
\begin{align}
\partial_xv-\partial_yu=-\frac{v}{2q^2}\partial_x\kappa+\frac{u}{2q^2}\partial_y\kappa-\frac{v}{\rho q^2}X+\frac{u}{\rho q^2}Y,
\label{align:ch18-18715}
\end{align}
将分别与\eqref{align:ch18-1821}$_2$ 与\eqref{align:ch18-18212}对照。

注意到变换 $(\rho,u,v)\mapsto(L,M,N)$ 可逆：
\begin{align}
u^2=-p(N-p),\qquad v^2=-p(L-p),\qquad M^2=(N-p)(L-p).
\label{align:ch18-18716}
\end{align}
事实上，\eqref{align:ch18-18716}中第三式确定 $p$，从而也确定 $\rho=-1/p$ 为 $(L,M,N)$ 的函数；随后\eqref{align:ch18-18716}中前两式同样把 $u$、$v$ 表示为 $(L,M,N)$ 的函数。

状态方程为\eqref{align:ch18-18711}的气体，其声速平方为 $c^2(\rho)=1/\rho^2$，故流动在 $\kappa>0$ 处亚声速，在 $\kappa<0$ 处超声速。由\eqref{align:ch18-18711}与\eqref{align:ch18-18712}，
\begin{align}
\rho=\frac{1}{\sqrt{q^2+\kappa}},\qquad p=-\sqrt{q^2+\kappa}.
\label{align:ch18-18717}
\end{align}

下文我们默认：$\rho$ 与 $p$ 作为 $q$ 的函数，已代入\eqref{align:ch18-1879}、\eqref{align:ch18-18714}与\eqref{align:ch18-18715}，于是这些方程只含两个未知量 $u$ 与 $v$。

在正则等距浸入的范围内，\eqref{align:ch18-1879}等价于\eqref{align:ch18-18714}、\eqref{align:ch18-18715}，二者中任一方程组都可用来确定浸入。然而对奇异等距浸入情形并非如此：此时\eqref{align:ch18-1879}仍成立，但\eqref{align:ch18-18714}与\eqref{align:ch18-18715}未必有效。上述内容应与\cref{section:ch18-定常平面等熵气流}中等熵定常无旋平面气流的设定对照：在那里，光滑流动同时满足\eqref{align:ch18-1821}、\eqref{align:ch18-18212}与\eqref{align:ch18-18214}；而按通行做法，对含激波的流动保留\eqref{align:ch18-1821}$_1$、\eqref{align:ch18-18212}与\eqref{align:ch18-18214}，弃去\eqref{align:ch18-1821}$_{2,3}$。

方程组\eqref{align:ch18-1879}——及其等价方程组\eqref{align:ch18-18714}、\eqref{align:ch18-18715}——在流动亚声速即 $\kappa>0$ 时为椭圆型，在流动超声速即 $\kappa<0$ 时为双曲型，当 Gauss 曲率在 $\Omega$ 内变号时为椭圆--双曲混合型。

此处我们聚焦双曲情形，以后总设 $\kappa<0$；特别地，不失一般性设 $\kappa=-1$（在 $\kappa$ 非常数的一般情形，$L,M,N$ 须重新标度为 $L/\sqrt{-\kappa}$、$M/\sqrt{-\kappa}$ 与 $N/\sqrt{-\kappa}$）。

下面概述一个构造奇异等距浸入的方案，目标情形为 $\Omega=\{(x,y):-\infty<x<\infty,\ 0<y<\infty\}$。细节见\cref{section:ch18-注记}所引文献。

把 $y$ 视为“时间”变量，任务是对双曲方程组\eqref{align:ch18-1879}的 Cauchy 问题建立 $L^\infty(\Omega)$ 弱解 $(u,v)$ 的存在性，并在 $x$ 轴上给定初值。为此，我们采用消失黏性方法，考察抛物方程组
\begin{align}
\begin{cases}
\partial_x(\rho u^2+p)+\partial_y(\rho uv)=\varepsilon\partial_y^2(\rho u)+X\\
\partial_x(\rho uv)+\partial_y(\rho v^2+p)=\varepsilon\partial_y^2(\rho v)+Y,
\end{cases}
\label{align:ch18-18718}
\end{align}
其中 $\varepsilon>0$ 最终趋于零。目标是证明：在给定初值下，\eqref{align:ch18-18718}的 Cauchy 问题在 $\Omega$ 上有经典解 $(u_\varepsilon,v_\varepsilon)$，它在 $L^\infty(\Omega)$ 中对 $\varepsilon>0$ 一致有界。一旦实现这一点，$(u_\varepsilon,v_\varepsilon)$ 便通过\eqref{align:ch18-1877}与\eqref{align:ch18-1878}生成函数 $(L_\varepsilon,M_\varepsilon,N_\varepsilon)$，后者同样在 $L^\infty(\Omega)$ 中对 $\varepsilon>0$ 一致有界。于是可抽取序列 $\{\varepsilon_n\}$，$\varepsilon_n\to0$（$n\to\infty$），使 $\{(L_{\varepsilon_n},M_{\varepsilon_n},N_{\varepsilon_n})\}$ 在 $L^\infty(\Omega)$ 弱$^*$ 收敛到 $\Omega$ 上的函数 $(L,M,N)$。借助 div--curl 引理，重演前面确立方程组\eqref{align:ch18-1876}、\eqref{align:ch18-1877}解在 $L^\infty$ 弱$^*$ 下序列紧性的步骤，易证 $(L,M,N)$ 满足\eqref{align:ch18-1876}、\eqref{align:ch18-1877}，从而诱导出所期望的流形到 $\mathbb{R}^3$ 的奇异等距浸入。最后，为把等距浸入与气流的类比贯彻到底，再由 $(L,M,N)$ 经\eqref{align:ch18-18716}导出 $L^\infty(\Omega)$ 函数 $(u,v)$，它求解方程组\eqref{align:ch18-1879}的 Cauchy 问题。

显然，上述方案的成败取决于：抛物方程组\eqref{align:ch18-18718}的 Cauchy 问题解在 $\varepsilon>0$ 下一致有界。事实证明，在等价方程组的框架下处理这一问题更为高效，
\begin{align}
\partial_x(\rho u)+\partial_y(\rho v)=\varepsilon\frac{u}{q^2}\partial_y^2(\rho u)+\varepsilon\frac{v}{q^2}\partial_y^2(\rho v)+\frac{u}{q^2}X+\frac{v}{q^2}Y,
\label{align:ch18-18719}
\end{align}
\begin{align}
\partial_xv-\partial_yu=-\varepsilon\frac{v}{\rho q^2}\partial_y^2(\rho u)+\varepsilon\frac{u}{\rho q^2}\partial_y^2(\rho v)-\frac{v}{\rho q^2}X+\frac{u}{\rho q^2}Y,
\label{align:ch18-18720}
\end{align}
它由\eqref{align:ch18-18718}重演前面从\eqref{align:ch18-1879}导出\eqref{align:ch18-18714}、\eqref{align:ch18-18715}的步骤得到。构造在 $\varepsilon>0$ 下一致有界的解，通过为该方程组定位一个与 $\varepsilon$ 无关的有界不变区域、并选取落于其中的初值来实现。

期望抛物方程组\eqref{align:ch18-18719}、\eqref{align:ch18-18720}的不变区域由双曲母方程组\eqref{align:ch18-18714}、\eqref{align:ch18-18715}的 Riemann 不变量的等值线为边界。经冗长计算，在极坐标 $u=q\cos\theta$、$v=q\sin\theta$ 下可证
\begin{align}
z=\theta-\arccos\frac1q,\qquad w=\theta+\arccos\frac1q
\label{align:ch18-18721}
\end{align}
是\eqref{align:ch18-18714}、\eqref{align:ch18-18715}的 Riemann 不变量。特别地，
\begin{align}
\begin{cases}
u\cos z+v\sin z=1\\
u\cos w+v\sin w=1,
\end{cases}
\label{align:ch18-18722}
\end{align}
它表明 $z$ 与 $w$ 的等值线在 $u$-$v$ 平面上均为直线，故方程组\eqref{align:ch18-18719}、\eqref{align:ch18-18720}的预备不变区域应为菱形。事实上，在对度量 $[g_{ij}]$ 的某些假设下（例如适用于悬链面），具有上述规格的不变区域的技术性构造，已在\cref{section:ch18-注记}所引文献中完成。

\end{document}
